\section{The Bristle and the Thread: Holy Fear as a Motion of Grace}

A serious believer sometimes hesitates before the fear of hell. Has not St.\ John written that ``fear is not in charity: but perfect charity casteth out fear, because fear hath pain'' (1 John 4:18)? If we ought to love God for His own infinite goodness, is it not base or mercenary to shrink from offending Him because of Gehenna?

Preaching on that very verse in his ninth Homily on the First Epistle of St.\ John, St.\ Augustine answers with two homely images that unlock the entire Catholic doctrine of fear. The first comes from a cobbler's bench:

\begin{quote}
Then let fear make the beginning, because ``the fear of the Lord is the beginning of wisdom.'' Fear, so to say, prepares a place for charity. But when once charity has begun to inhabit, the fear which prepared the place for it is cast out. \dots\ As we see in sewing, the thread is introduced by means of the bristle; the bristle first enters, but except it come out the thread does not come into its place: so fear first occupies the mind, but the fear does not remain there, because it enters only in order to introduce charity.\endnote{St.\ Augustine, \textit{Homilies on the First Epistle of John} 9.4, NPNF1 vol.\ 7.}
\end{quote}

The second image comes from a surgeon's table. When rottenness festers inside a wound, the surgeon's knife makes the flesh smart more sharply while he cuts than it smarted while the poison was left undisturbed; yet ``it smarts more under the healing operation, but only that it may never smart when the healing is effected.'' So the fear of God cuts into a complacent conscience in order to clear the way for charity to heal the wound completely. ``Needs then must fear first enter in, that by it charity may come. Fear is the healing operation: charity, the sound condition.''

Then Augustine confronts an apparent contradiction in Scripture itself. St.\ John writes that perfect charity casts out fear (1 John 4:18), yet Psalm 18:10 sings that ``the fear of the Lord is chaste, enduring for ever.'' How can one fear be cast out while another fear endures forever? One Spirit fills two pipes of a double flute, Augustine answers, and their notes make harmony once we distinguish two kinds of fear. To show the difference, he sets two married women side by side:

\begin{quote}
One cannot better explain the difference between these two fears, the one which charity casteth out, the other chaste, which endureth for ever, than by putting the case of two married women, one of whom, you may suppose, is willing to commit adultery, delights in wickedness, only fears lest she be condemned by her husband. \dots\ Put the case that the other loves her husband, that she feels that she owes him chaste embraces, that she stains herself with no uncleanness of adultery; she wishes for the presence of her husband. \dots\ Question the one, Dost thou fear thine husband? she answers, I do. Question the other, whether she fears her husband; she answers, I do fear him. The voice is one, the mind diverse. Now then let them be questioned, Why? The one saith, I fear my husband, lest he should come: the other saith, I fear my husband, lest he depart from me. The one saith, I fear to be condemned: the other, I fear to be forsaken.\endnote{St.\ Augustine, \textit{Homilies on the First Epistle of John} 9.5--6, NPNF1 vol.\ 7.}
\end{quote}

Few paragraphs in Christian theology reach deeper into the human heart. Both wives fear; one dreads the judge's arrival, the other dreads the beloved's departure. The fear that shrinks from being condemned to hell is the cobbler's bristle: imperfect, yet indispensable when a sinner is still entangled in sin. The chaste fear that dreads being forsaken by God is born of love itself and abides with the saints in eternity.

St.\ Thomas Aquinas turns Augustine's pastoral distinction into a crystal-clear map in the \textit{Secunda Secundae} (II-II, q.\ 19). First, Aquinas shows that God is the object of both hope and fear without any contradiction: ``in us the consideration of His justice gives rise to fear, but the consideration of His mercy gives rise to hope, so that, accordingly, God is the object of both hope and fear, but under different aspects'' (a.\ 1, ad 2).\endnote{St.\ Thomas Aquinas, \textit{Summa Theologiae} II-II, q.\ 19, a.\ 1, ad 2, trans.\ Fathers of the English Dominican Province.} Next, he distinguishes four kinds of fear: \textit{worldly fear}, which shrinks from temporal loss so much that it turns away from God and is always evil; \textit{servile fear}, which shrinks from the evil of punishment inflicted by God; \textit{initial fear}, in which a beginner moving toward grace fears both punishment and separation from God; and \textit{filial} or \textit{chaste fear}, a gift of the Holy Ghost, which shrinks above all from the evil of fault that would separate the soul from its Father (aa.\ 2, 4--8).

Is servile fear---the fear of hell---good or evil? Aquinas answers with exact precision in Article 4: ``servile fear is substantially good, but its servility is evil.''\endnote{St.\ Thomas Aquinas, \textit{Summa Theologiae} II-II, q.\ 19, a.\ 4.} What is the \textit{servility}? That twisted attachment of the slave who still loves sin as his end and would gladly commit it if only there were no fire. What is the \textit{substance} of servile fear? Simply shrinking from eternal punishment as an evil contrary to one's true good, without placing one's end in sin. In that substance, fear of hell is good, remains compatible with charity as a subordinate motive (a.\ 6), and serves as ``the beginning of wisdom'' by disposing a man from without---``in so far as he refrains from sin through fear of punishment, and is thus fashioned for the effect of wisdom'' (a.\ 7).\endnote{St.\ Thomas Aquinas, \textit{Summa Theologiae} II-II, q.\ 19, aa.\ 6--7.}

Three centuries later, the Council of Trent raised this patristic and Thomistic insight into solemn conciliar teaching against the claim that fear of hell makes the sinner a hypocrite. In Session XIV, Chapter 4, on the Sacrament of Penance, the council declares:

\begin{quote}
And as to that imperfect contrition, which is called attrition, because that it is commonly conceived either from the consideration of the turpitude of sin, or from the fear of hell and of punishment, It declares that if, with the hope of pardon, it exclude the wish to sin, it not only does not make a man a hypocrite, and a greater sinner, but that it is even a gift of God, and an impulse of the Holy Ghost,---who does not indeed as yet dwell in the penitent, but only moves him,---whereby the penitent being assisted prepares a way for himself unto justice.\endnote{Council of Trent, Session XIV (25 November 1551), Doctrine on the Sacrament of Penance, ch.\ 4, trans.\ J.\ Waterworth (London: Dolman, 1848); DS 1678.}
\end{quote}

And in Canon 5 of the same session, Trent anathematizes anyone who says that sorrow arising from examining and detesting one's sins, pondering their gravity, ``the loss of eternal blessedness, and the eternal damnation which he has incurred, with the purpose of a better life,'' is not a true and profitable sorrow or does not prepare for grace.\endnote{Council of Trent, Session XIV, Canon 5, trans.\ Waterworth; DS 1705.} Imperfect contrition does not justify a sinner apart from the Sacrament of Penance, yet inside the Sacrament it suffices for valid absolution, and even before Confession it is already the Holy Ghost moving the soul from outside to break its covenant with death.

This is why Scripture pairs its sternest warnings with the tenderest assurances of God's saving will. God swears by His own life through the prophet Ezechiel: ``As I live, saith the Lord God, I desire not the death of the wicked, but that the wicked turn from his way, and live. Turn ye, turn ye from your evil ways: and why will you die, O house of Israel?'' (Ezech 33:11; cf.\ 18:23, 32). St.\ Paul teaches that God our Saviour ``will have all men to be saved and to come to the knowledge of the truth'' (1 Tim 2:4), and St.\ Peter declares that the Lord deals patiently for our sake, ``not willing that any should perish, but that all should return to penance'' (2 Pet 3:9). Most decisively of all, ``God so loved the world, as to give his only begotten Son: that whosoever believeth in him may not perish, but may have life everlasting. For God sent not his Son into the world, to judge the world: but that the world may be saved by him'' (John 3:16--17).

Holy fear never stands alone, and it never leads to despair. Despair forgets the Cross; presumption forgets the Judge. Between those two ditches runs the narrow royal road that St.\ Paul sets before the Philippians: ``with fear and trembling work out your salvation. For it is God who worketh in you, both to will and to accomplish, according to his good will'' (Phil 2:12--13). Even the Apostle of the Gentiles wrote, ``I chastise my body and bring it into subjection: lest perhaps, when I have preached to others, I myself should become a castaway'' (1 Cor 9:27). When fear of hell drives a sinner to the confessional, breaks off a habit of sin, awakens mercy toward the poor, and deepens gratitude for the Blood of Christ, the bristle has done its work and drawn the thread of charity into the heart.
